% Figure: the partial-sum reconstruction of a square wave from its odd
% harmonics, showing the Gibbs overshoot at the discontinuity, and the
% harmonic-amplitude spectrum. pgfplots. The square wave has period 2 and
% unit amplitude; harmonics fall as 1/n for odd n only.
\begin{figure}[htbp]
\centering
% === left: partial sums in the time domain ===
\begin{tikzpicture}
\begin{axis}[
  width=7.2cm, height=5.4cm,
  xlabel={$t/T$}, ylabel={$v(t)$},
  xmin=0, xmax=1, ymin=-1.5, ymax=1.5,
  axis lines=left,
  tick label style={font=\tiny},
  label style={font=\scriptsize},
  legend style={font=\tiny, at={(0.98,0.02)}, anchor=south east, draw=none, fill=none},
  clip=false,
]
% one term: (4/pi) sin(2 pi t)
\addplot[enggray, thick, domain=0:1, samples=200]
  {(4/pi)*sin(deg(2*pi*x))};
\addlegendentry{$N=1$}
% three terms (n = 1,3,5): add (4/pi)[ sin/1 + sin3/3 + sin5/5 ]
\addplot[engorange, thick, domain=0:1, samples=300]
  {(4/pi)*( sin(deg(2*pi*x)) + sin(deg(6*pi*x))/3 + sin(deg(10*pi*x))/5 )};
\addlegendentry{$N=5$}
% many terms (n up to 19): full Gibbs ear visible
\addplot[engblue, very thick, domain=0:1, samples=600]
  {(4/pi)*( sin(deg(2*pi*x)) + sin(deg(6*pi*x))/3 + sin(deg(10*pi*x))/5
   + sin(deg(14*pi*x))/7 + sin(deg(18*pi*x))/9 + sin(deg(22*pi*x))/11
   + sin(deg(26*pi*x))/13 + sin(deg(30*pi*x))/15 + sin(deg(34*pi*x))/17
   + sin(deg(38*pi*x))/19 )};
\addlegendentry{$N=19$}
\draw[engred, dashed] (axis cs:0,1) -- (axis cs:0.5,1);
\draw[engred, dashed] (axis cs:0.5,-1) -- (axis cs:1,-1);
\end{axis}
\end{tikzpicture}
\hspace{0.4cm}
% === right: harmonic amplitude spectrum ===
\begin{tikzpicture}
\begin{axis}[
  width=6.2cm, height=5.4cm,
  xlabel={harmonic $n$}, ylabel={amplitude $|c_n|$},
  xmin=0, xmax=12, ymin=0, ymax=1.4,
  axis lines=left,
  tick label style={font=\tiny},
  label style={font=\scriptsize},
  ybar, bar width=3pt,
  xtick={1,3,5,7,9,11},
  clip=false,
]
% (4/pi)/n for odd n; even harmonics are zero
\addplot[engblue, fill=engblue] coordinates {
  (1,1.2732) (3,0.4244) (5,0.2546) (7,0.1819) (9,0.1415) (11,0.1157)};
\end{axis}
\end{tikzpicture}
\caption[Fourier reconstruction and spectrum of a square wave]{A unit
square wave reconstructed from its Fourier series (left) and its
harmonic-amplitude spectrum (right). The square wave contains only odd
harmonics, whose amplitudes fall as $4/(\pi n)$, drawn as the bar
spectrum. Adding more terms sharpens the edges, but the partial sum
keeps a fixed overshoot of about nine percent at the discontinuity, the
Gibbs phenomenon, which narrows but never shrinks as the harmonic count
rises. The $1/n$ amplitude roll-off is why the power in the higher
harmonics falls as $1/n^2$ and the total harmonic content of a
fast-edged waveform concentrates in its first few harmonics.}
\label{fig:vol04ch09:fourier-square}
\end{figure}
